\documentclass[10pt,a4paper]{article}

\usepackage{fontspec}
\setmainfont{PT Sans}
\usepackage[a4paper,margin=18mm]{geometry}
\usepackage{tabularx}
\usepackage{booktabs}
\usepackage{xcolor}
\usepackage{enumitem}
\usepackage{amsmath}
\usepackage{hyperref}

\hypersetup{colorlinks=true,linkcolor=black,urlcolor=blue}
\setlength{\parindent}{0pt}
\setlength{\parskip}{5pt}
\setlist[itemize]{leftmargin=*,topsep=2pt,itemsep=2pt}

\definecolor{headerblue}{HTML}{1F4E79}
\definecolor{lightgray}{HTML}{F3F5F7}

\begin{document}

{\Large\bfseries Отчет по аугментации MRI train pairs}

\vspace{4pt}
{\color{headerblue}\rule{\linewidth}{1.2pt}}

\section*{Кратко}

На сервере сгенерирован расширенный train dataset для задачи T1 post-contrast
query к T2 target retrieval. Аугментация применяется к положительным парам из
\texttt{dataset1/train\_pairs.csv}. Для каждой augmented copy T1 и T2
преобразуются независимо, чтобы сохранить same-subject positive pair, но
убрать идеальное геометрическое совпадение между модальностями.

\section*{Что сгенерировано}

\begin{tabularx}{\linewidth}{@{}lX@{}}
\toprule
Пункт & Значение \\
\midrule
Оригинальные train pairs & 350 \\
Augmented pairs всего & 1750 \\
Итоговых train pairs в CSV & 2100 \\
Строк в CSV с header & 2101 \\
Augmented NIfTI files всего & 3500 \\
Старые augmented files & 1000 \\
Новые stronger-extra files & 2500 \\
Проверка путей в CSV & missing paths = 0 \\
Проверка pair id & unique pair ids = 2100 \\
\bottomrule
\end{tabularx}

\section*{Где лежит}

\begin{itemize}
  \item Финальный CSV:
  \path{/app/ehl-paris-medical-image-retrieval/dataset1/train_pairs_aug_geom_contrast_1k_plus_stronger.csv}
  \item Старые augmented files:
  \path{/app/ehl-paris-medical-image-retrieval/dataset1_aug_geom_contrast_1k}
  \item Новые stronger-extra augmented files:
  \path{/app/ehl-paris-medical-image-retrieval/dataset1_aug_geom_contrast_stronger_extra}
  \item Логи генерации:
  \path{/app/runs/generate_aug_geom_contrast_stronger_extra}
\end{itemize}

\section*{Использованные наборы аугментаций}

\begin{tabularx}{\linewidth}{@{}lccX@{}}
\toprule
Набор & Пар & Copies & Что добавлено \\
\midrule
\texttt{geom\_contrast} & 250 & 2 & Старый набор: \texttt{aug00-aug01} для первых 250 пар. \\
\texttt{geom\_contrast\_stronger} & 250 & 3 & Новые \texttt{aug02-aug04} для первых 250 пар. \\
\texttt{geom\_contrast\_stronger} & 100 & 5 & Новые \texttt{aug00-aug04} для оставшихся 100 пар. \\
\bottomrule
\end{tabularx}

\section*{Трансформации и значения}

\begin{tabularx}{\linewidth}{@{}lccX@{}}
\toprule
Трансформация & \texttt{geom\_contrast} & \texttt{stronger} & Что делает \\
\midrule
Rotation & $\pm18^\circ$ & $\pm22^\circ$ & Поворачивает 3D volume вокруг осей x/y/z. \\
Translation & $\pm12$ voxels & $\pm16$ voxels & Сдвигает volume в пространстве. \\
Scale & 0.90-1.10 & 0.85-1.15 & Немного растягивает или сжимает volume. \\
Elastic deformation & p=0.45, 3-8 voxels & p=0.60, 4-10 voxels & Smooth nonlinear deformation: разные области смещаются по-разному. \\
Bias field & p=0.30 & p=0.40 & Плавно меняет яркость по объему, имитируя MRI inhomogeneity. \\
Gamma contrast & p=0.30, gamma 0.75-1.35 & p=0.40, gamma 0.75-1.35 & Нелинейно меняет контраст. \\
Intensity scale/shift & p=0.30 & p=0.40 & Линейно меняет яркость: scale плюс небольшой offset. \\
Gaussian noise & p=0.30 & p=0.40 & Добавляет небольшой случайный шум. \\
\bottomrule
\end{tabularx}

\section*{Итоговая логика}

Финальный CSV содержит 350 оригинальных пар и 1750 аугментированных пар.
Для первых 250 исходных пар теперь есть 5 augmented copies. Для оставшихся
100 исходных пар также есть 5 augmented copies. Всего это дает:

\[
350 \times 5 = 1750 \text{ augmented pairs}
\]

\[
1750 \times 2 = 3500 \text{ augmented NIfTI files}
\]

Обучение модели и Kaggle submit в рамках этой генерации не запускались.

\end{document}
